\documentclass[10pt,a4paper]{amsart}
\usepackage[margin=19mm]{geometry}
\usepackage{amsmath,amssymb,mathtools}
\usepackage[scr=rsfs,cal=euler]{mathalfa}
\usepackage{xfrac}
\usepackage[unicode,hidelinks,bookmarks=false]{hyperref}
\newcommand{\ie}{i.e.\ }
\newcommand{\cf}{cf.\ }
\newcommand{\resp}{respectively\ }
\newcommand{\Subsection}{\subsection}
\providecommand{\nobreakdash}{\nobreak\hbox{-}}
\setlength{\emergencystretch}{2em}
\multlinegap0pt
\usepackage{xcolor}
\usepackage{tcolorbox}
\usepackage{tikz}
\newtheorem{lemma}{Lemma}
\title{Effective Context in Transformers: An Analysis of Fragmentation and Tokenization}
\author{Amirmehdi Jafari Fesharaki, Mohammadamin Rami, Aslan Tchamkerten}
\date{Source: arXiv:2605.13485v1 (CC BY 4.0)}
\begin{document}
\maketitle

\noindent\emph{Source and licence.} Effective Context in Transformers: An Analysis of Fragmentation and Tokenization. Version arXiv:2605.13485v1 (CC BY 4.0), distributed by arXiv under the Creative Commons Attribution 4.0 International licence, http://creativecommons.org/licenses/by/4.0/, as recorded in the arXiv licence metadata for this exact version and shown on the abstract page https://arxiv.org/abs/2605.13485v1. Full licence notice: \texttt{licenses/arxiv-2605.13485-CC-BY-4.0.txt}.\par\smallskip

\noindent\textit{Editorial scope.} Counted unit: the article's lemma lem:hh-typical-source-span (source0.tex lines 2467--2494). The article's complete original proof of it is delivered with it (source0.tex lines 2496--2568, one \texttt{proof} environment of 73 lines). No mathematics is written, repaired or re-derived: the statement and the proof are the article's own lines, in the article's own notation, and no retained line is substituted or re-flowed.  Retained uncounted as the source-local context the proof needs: lines 2412--2427.  Nothing was omitted from any retained span, so the package carries no omission record.  No retained line contains a citation, so the wrapper carries no bibliography and no bibliography entry.  No retained line contains a figure environment, an image-inclusion command or drawing code, so no image is delivered and the package declares no asset dependency.  Editorial formatting only, in addition: an AMS \texttt{amsart} preamble that restates the article's own declarations and macros used by the retained lines, namely the environments \texttt{lemma} on separate counters, together with the standard packages those lines need.  The retained environments print in this document as Lemma 1 in order of appearance, whereas the article prints the same items with the numbering of its own full document.  The complete native source of the article as distributed in the arXiv e-print for this exact version is bundled unchanged as \texttt{sources/200507/source0.tex}.
    \begin{lemma}[Typical heavy-hitting implies long tokens with high probability]
    \label{lem:hh-long-tokens}
    Assume $\mathcal Z$ is $(\beta,\eta)$-typically heavy-hitting. Define
    \[
        \ell_d
        :=
        \frac{\beta\log_2 d}{\log_2(1/\delta)}.
    \]
    Then
    \[
        \Pr_T\!\left[
            |\Gamma_{\mathcal Z}(T)|<\ell_d
        \right]
        \le \eta .
    \]
    \end{lemma}

    \begin{lemma}[Typical heavy-hitting implies typical source span]
    \label{lem:hh-typical-source-span}
    Assume $\mathcal Z$ is $(\beta,\eta)$-typically heavy-hitting. Fix a
    token-window length $w$, and define
    \[
        \ell_d
        :=
        \frac{\beta\log_2 d}{\log_2(1/\delta)}.
    \]
    Let
    \[
        w_d
        :=
        \left\lfloor \frac{3}{4}w\ell_d \right\rfloor .
    \]
    Then
    \[
        \Pr\!\left[
            S(Z_{i-w}^{i-1}) < w_d
        \right]
        \le
        4\eta.
    \]
    Equivalently, $\Pi_{\mathcal Z}$ is
    \[
        \left(4\eta,\; w,\; w_d\right)\text{-typical}.
    \]
    \end{lemma}

    \begin{proof}
    For each token $Z_j$, define the bad-token indicator
    \[
        B_j
        :=
        \mathbf 1\left\{
            |\Gamma_{\mathcal Z}(Z_j)|<\ell_d
        \right\}.
    \]
    By Lemma~\ref{lem:hh-long-tokens},
    \[
        \mathbb E[B_j]=\Pr(B_j=1)\le \eta.
    \]
    For the window $Z_{i-w}^{i-1}$, let
    \[
        B
        :=
        \sum_{j=i-w}^{i-1}B_j
    \]
    be the number of bad tokens in the window. By linearity of expectation,
    \[
        \mathbb E[B]\le w\eta.
    \]
    
    Every good token has source length at least $\ell_d$, and every bad token has
    source length at least $1$. Therefore,
    \[
        S(Z_{i-w}^{i-1})
        =
        \sum_{j=i-w}^{i-1}|\Gamma_{\mathcal Z}(Z_j)|
        \ge
        (w-B)\ell_d + B.
    \]
    In particular, if $B\le w/4$, then
    \[
        S(Z_{i-w}^{i-1})
        \ge
        \left(w-\frac{w}{4}\right)\ell_d
        =
        \frac{3}{4}w\ell_d
        \ge
        w_d.
    \]
    Thus,
    \[
        \left\{
            S(Z_{i-w}^{i-1}) < w_d
        \right\}
        \subseteq
        \left\{
            B>\frac{w}{4}
        \right\}.
    \]
    By Markov's inequality,
    \[
        \Pr\!\left[
            B>\frac{w}{4}
        \right]
        \le
        \frac{\mathbb E[B]}{w/4}
        \le
        4\eta.
    \]
    Therefore,
    \[
        \Pr\!\left[
            S(Z_{i-w}^{i-1}) < w_d
        \right]
        \le
        4\eta.
    \]
    This is exactly the $\left(4\eta,w,w_d\right)$-typicality condition.
    \end{proof}

\end{document}
